% !TeX root = ../../Book3.tex
% 纲要标题：### 18.3 曲面源上的小能量正则性
% 纲要位置：几何学计划纲要.md，第 2836 行。
% 正文状态：未撰写；本文件目前是框架。

% 纲要内容开始
% 
% 完整证明的主版本固定为二维圆盘内的光滑调和映射、光滑紧黎曼目标及内域估计。先对光滑解族建立一致界，不把所有高维弱解都归入同一小能量定理。
% 
% * 二维能量缩放与临界性
% * 光滑调和映射的 Bochner 微分不等式
% * 小能量阈值、均值估计和梯度界
% 
%   证明存在由目标几何决定的 $\varepsilon_0,C>0$，使 $E(u;B_r)<\varepsilon_0$ 时 $\sup_{B_{r/2}}|du|^2\le Cr^{-2}E(u;B_r)$；在曲面坐标中追踪背景度量的控制常数。
% 
% * 内域高阶估计及光滑解族的紧性
% * 能量集中点、重缩放及奇点形成
% * 一般弱解与高维部分正则性的独立范围
% 
%   二维一般弱调和映射的正则性及高维极小、平稳映射的部分正则性分别列出外部定理的准确对象和假设；单调性公式仅用于已经验证相应驻值条件的映射。
% 
% 纲要内容结束

\section{曲面源上的小能量正则性}
\label{g3:sec:1803}

待撰写

\subsection{二维能量缩放与临界性}
\label{g3:subsec:1803:01}
待撰写

\subsection{光滑调和映射的 Bochner 微分不等式}
\label{g3:subsec:1803:02}
待撰写

\subsection{小能量阈值、均值估计和梯度界}
\label{g3:subsec:1803:03}
待撰写

\subsection{内域高阶估计及光滑解族的紧性}
\label{g3:subsec:1803:04}
待撰写

\subsection{能量集中点、重缩放及奇点形成}
\label{g3:subsec:1803:05}
待撰写

\subsection{一般弱解与高维部分正则性的独立范围}
\label{g3:subsec:1803:06}
待撰写
% 纲要条目排为编号小节标题，并进入目录与书签；说明段仍保留于纲要注释中。
